\section{Introduction}
\label{sec:intro}

The growth of audiovisual content has driven the development of methods capable of analysing, indexing and automatically retrieving visual information. Within this landscape, multimodal video--text retrieval has consolidated as a relevant area at the intersection of computer vision and natural language processing. An especially challenging scenario is that of \emph{egocentric video}, captured from a first-person perspective using head- or body-mounted cameras. Such recordings register natural interactions between hands, objects and the surrounding environment, introducing strong viewpoint, illumination and motion variability. Models must therefore interpret actions and contextual relationships at fine granularity in order to establish precise correspondences between visual content and textual descriptions.

The reference resource for this problem is the \href{https://epic-kitchens.github.io/}{EPIC-KITCHENS-100} dataset~\cite{damen2022rescaling}, which contains over 100\,h of egocentric kitchen video and approximately 90{,}000 narrated action segments. The associated \href{https://www.codabench.org/competitions/12008/\#/pages-tab}{Multi-Instance Retrieval (MIR) challenge} formalises the task: given a textual query, the system must rank 9{,}668 trimmed video segments from the test split (and vice versa) by their semantic relevance, where ground-truth relevance is provided as a continuous \emph{soft-label matrix} rather than as one-to-one binary correspondences~\cite{damen2022rescaling,wray2021semantic}. This formulation acknowledges that several segments may match a query at different degrees of similarity and so, is evaluated with order-sensitive metrics like \emph{normalized discounted cumulative gain} (nDCG)~\cite{jarvelin2002ndcg} and \emph{mean average precision} (mAP) rather than with the conventional Recall@K used in third-person retrieval. Despite recent progress in multimodal foundation models, semantic alignment between egocentric video and natural language remains difficult due to dynamic perspective changes, occlusions and multi-object interactions hinder the extraction of robust visual representations, and the fact that most existing losses do not exploit the graded relevance signal provided by the challenge.

% \paragraph{Contributions.}
% This work systematically explores the EK-100 MIR task through three progressively stronger approaches that share an iterative development methodology grounded on the same evaluation protocol. Concretely, we contribute:

% \begin{enumerate}
% \setlength{\itemsep}{2pt}
%   \item A reproduction and analysis of a \textbf{JPoSE-based zero-fine-tuning baseline}~\cite{wray2019jpose}, in which verb, noun and action sub-embeddings of pre-trained Part-of-Speech encoders are concatenated to produce a $9668\times 3842$ similarity matrix. This baseline reaches \textbf{53.53\,nDCG} on the official Codabench leaderboard and provides the lower bound of our study.

%   \item A \textbf{score-level ensemble} that fuses three pre-trained retrieval models (JPoSE, MI-MM~\cite{miech2020milnce} and MMEN~\cite{wray2019jpose}) through row-wise min-max normalization followed by equal-weight averaging, optionally followed by a \textbf{graph-diffusion re-ranking} step that propagates similarity scores through visual and textual k-nearest-neighbour graphs in the spirit of~\cite{iscen2017diffusion}. Re-ranking provides a further uplift over the ensemble alone, reaching \textbf{55.82\,nDCG}.

%   \item A complete \textbf{end-to-end fine-tuning pipeline} that trains an AVION ViT-L/14 dual encoder~\cite{zhao2023avion} with the Symmetric Multi-Similarity (SMS) loss~\cite{wang2024symmetricmultisimilaritylossepickitchens100}, which consumes the soft-label relevancy matrices directly during training, complemented by a \emph{horizontal-flip + 32-frame} test-time augmentation strategy. The final iteration, trained for sixteen epochs on a single GPU, attains \textbf{69.68\,nDCG} on the official Codabench leaderboard---a \textbf{+16.15} absolute gain over the JPoSE baseline.

%   \item A discussion of where each gain originates, whether encoder capacity, soft-label supervision, test-time augmentation or additional epochs, together with a report of the primary limitations encountered (single-GPU compute budget and unavailability of the test-set relevancy matrix for offline evaluation) and a future-work proposal centered on a \emph{model-soup}-style ensemble~\cite{wortsman2022modelsoups} of differently-trained ViT-L+SMS instances.
% \end{enumerate}
